\sectionguard
\section{Reception and Ecumenical Boundaries}

Two of the four dogmas are held, in substance, by Christians who are not
in communion with Rome; two are contested, and contested for reasons that
are not primarily about Mary. A reference that states the Catholic
doctrine owes its reader an accurate map of where the disagreement
actually lies. This chapter draws it from the other communions' own
documents wherever such documents could be identified, and marks the
level of each.

\subsection*{The Orthodox East}

The Christological dogmas are common ground, and the shared inheritance
is not a modern convergence but the same councils. The Orthodox
Churches confess the \term{Theotokos} of Ephesus and Chalcedon, keep the
\term{aeiparthenos} of Constantinople II, celebrate the Dormition
(\term{Koimesis}) on 15 August, and honor Mary in a liturgical register
at least as rich as the Latin one. \work{Lumen gentium} 69 registers the
fact with evident relief: the Council rejoices that among the separated
brethren ``especially among the Orientals,'' many ``with devout mind and
fervent impulse give honor to the Mother of God, ever
virgin.''\endnote{LG 69, Vatican English text read at
\url{https://www.vatican.va/archive/hist_councils/ii_vatican_council/documents/vat-ii_const_19641121_lumen-gentium_en.html}.}

The difficulty is the Immaculate Conception, and the 1854 definition drew
an explicit answer. In 1895 the Ecumenical Patriarch Anthimus VII and his
synod replied to Leo XIII's \work{Praeclara gratulationis} with an
encyclical listing the innovations they charged against the Roman Church.
Article XIII of that list is the Marian one, and its wording repays exact
quotation:

\begin{studybox}{Patriarchal and Synodal Encyclical of Constantinople,
1895, art.~XIII (English witness; not a Catholic act)}
``The one holy, catholic and apostolic Church of the seven Ecumenical
Councils teaches that the supernatural incarnation of the only-begotten
Son and Word of God, of the Holy Ghost and the Virgin Mary, is alone pure
and immaculate; but the Papal Church scarcely forty years ago again made
an innovation by laying down a novel dogma concerning the immaculate
conception of the Mother of God and ever-Virgin Mary, which was unknown
to the ancient Church (and strongly opposed at different times even by
the more distinguished among the papal theologians).''\endnote{
Patriarchal Encyclical of 1895, art.~XIII, English text read at
\url{http://orthodoxinfo.com/ecumenism/encyc_1895.aspx}. The witness is a
private Orthodox website reproducing a nineteenth-century English
translation; neither the translator nor the underlying Greek edition is
identified on the page, and no collation was performed. The document is
the synodal reply of Anthimus VII to Leo XIII's \work{Praeclara
gratulationis} (1894); an earlier article of the same encyclical also
lists ``the immaculate conception of the ever-Virgin'' among the points
the Latins would have to prove from the Fathers and the councils.}
\end{studybox}

Three things in that sentence deserve notice, because each is a precise
datum rather than a general hostility. First, the objection is
\emph{historical}: the doctrine ``was unknown to the ancient Church.''
Second, it appeals to the Latin tradition's own internal
dissent---``strongly opposed at different times even by the more
distinguished among the papal theologians,'' which is the medieval debate
of Chapter 5 turned into an ecumenical argument. Third, the same sentence
that rejects the dogma calls Mary ``the Mother of God and ever-Virgin'':
the two ancient dogmas are affirmed in the act of denying the modern one.

Behind the historical objection lies a difference in the doctrine of
original sin itself. The Greek tradition's account of what is inherited
from Adam---mortality and corruption, with sin as its consequence---does
not pose the question to which ``preserved from the stain of original
sin at the first instant'' is the answer, so that the definition can
appear to Orthodox theologians both unnecessary and distorting. This
reference states that difference at the level of a widely reported
theological characterization and does not attempt to adjudicate
it.\endnote{Reported at that level: no Orthodox conciliar or synodal act
defining the doctrine of ancestral sin was identified and read for this
publication, and the characterization is therefore not attributed to any
particular authority. What is documented above is the 1895 synodal
encyclical's own reasoning.} On the Assumption the disagreement is of a
different kind: the content---the Dormition and the glorification of the
\term{Theotokos}---is held in the East's liturgy, and \work{MD} itself
quotes the Byzantine liturgy in evidence; what is rejected is the mode,
a dogma defined by the Bishop of Rome alone.

\subsection*{The Reformation traditions}

The Reformers did not begin as Marian minimalists, and the sixteenth
century's confessional documents show it. The \work{Smalcald Articles} of
1537, Luther's own, place among the articles ``concerning which there is
no contention or dispute, since we on both sides confess them'' the
statement that the Son ``was conceived, without the cooperation of man,
by the Holy Ghost, and was born of the pure, holy [and always] Virgin
Mary.''\endnote{\work{The Smalcald Articles} (1537), Part I, art.~IV,
in the translation of F. Bente and W. H. T. Dau, \work{Triglot
Concordia: The Symbolical Books of the Ev.\ Lutheran Church} (St.~Louis:
Concordia, 1921), pp.~453--529; public-domain text read at
\url{https://www.gutenberg.org/files/273/273-h/273-h.htm}. The square
brackets are the Triglot's own editorial convention marking wording
carried by one of the parallel originals; that convention, not this
reference, governs the bracketed words, and no claim about the German and
Latin recensions is made here beyond what the edition prints.} Both the
divine motherhood (the article is a Christological one) and the perpetual
virginity thus stand inside a foundational Lutheran confession, listed
among the points not in dispute with Rome.

The objection that developed is structural rather than Marian, and it
falls on the two modern definitions. It has two parts, and the
Anglican--Roman Catholic International Commission states both with
unusual clarity. The first is the sufficiency of Scripture, in the words
of Article VI of the Thirty-Nine Articles, which ARCIC quotes: ``Holy
Scripture containeth all things necessary to salvation: so that
whatsoever is not read therein, nor may be proved thereby, is not to be
required of any man, that it should be believed as an article of the
Faith.'' The question this raises, as ARCIC puts it, is ``whether these
doctrines concerning Mary are revealed by God in a way which must be held
by believers as a matter of faith.'' The second is authority: the
doctrines were defined by the Bishop of Rome ``independent of a
Council.''\endnote{\work{Mary: Grace and Hope in Christ} (the Seattle
Statement of ARCIC II, agreed 2004, published 2005), \S\S60, 62, read at
\url{https://www.christianunity.va/content/unitacristiani/en/dialoghi/sezione-occidentale/comunione-anglicana/dialogo/arcic-ii/en.html}.}

\subsection*{ARCIC's Seattle Statement (2005)}

The most developed ecumenical treatment of the two modern dogmas is the
Seattle Statement, and its structure is worth reporting exactly, because
it is easy to over- or under-read.

What it agrees. On the Assumption, ``given the understanding we have
reached concerning the place of Mary in the economy of hope and grace, we
can affirm together the teaching that God has taken the Blessed Virgin
Mary in the fullness of her person into his glory as consonant with
Scripture,'' noting that the dogma ``does not adopt a particular position
as to how Mary's life ended.'' On the Immaculate Conception, ``we can
affirm together that Christ's redeeming work reached `back' in Mary to
the depths of her being, and to her earliest beginnings.'' Summing both:
``the teaching about Mary in the two definitions of 1854 and 1950,
understood within the biblical pattern of the economy of grace and hope
outlined here, can be said to be consonant with the teaching of the
Scriptures and the ancient common traditions.''\endnote{Same witness,
\S\S58--60. The statement adds, of the Immaculate Conception, an
observation that is doctrinally acute: ``the negative notion of
`sinlessness' runs the risk of obscuring the fullness of Christ's saving
work. It is not so much that Mary lacks something which other human
beings `have', namely sin, but that the glorious grace of God filled her
life from the beginning.''}

What it does not resolve. ``Consonant with Scripture'' is not the same as
``revealed by God and therefore \term{de fide}''; ARCIC says so, and
leaves the question of required assent to a wider process of what it
calls re-reception, adding that Anglicans have asked whether acceptance
of the two definitions would be a condition of restored communion, and
that ``Roman Catholics find it hard to envisage a restoration of
communion in which acceptance of certain doctrines would be requisite for
some and not for others.''\endnote{Same witness, \S\S60--63. \S61 also
observes that the 1854 and 1950 phrases ``revealed by God'' and
``divinely revealed'' reflect the theology of revelation dominant when
they were made and are to be understood in the light of \work{Dei
verbum}---an interpretative move which is the Commission's, not a
magisterial act.}

What it is. The document says so itself, in a paragraph headed ``The
Status of the Document'': it ``is a joint statement of the Commission,''
whose publication the appointing authorities ``have allowed \ldots{} so
that it may be widely discussed,'' and it ``is not an authoritative
declaration by the Roman Catholic Church or by the Anglican Communion,
who will study and evaluate the document in due course.'' This reference
cites it accordingly: the most careful available statement of where the
difficulty lies, published on the Holy See's own Christian-unity site,
carrying exactly the authority of a dialogue text and no more.\endnote{
Same witness, ``The Status of the Document,'' preceding the statement;
the co-chairmen's preface is signed by Alexander J. Brunett and Peter F.
Carnley at Seattle on the Feast of the Presentation, 2 February 2004. The
Commission's own \S63 makes any ``re-reception'' conditional on ``a
mutual re-reception of an effective teaching authority in the Church.''}

\subsection*{Where the disagreement actually lies}

Gathering the evidence: the divine motherhood is not in dispute among the
communions treated here; the perpetual virginity is confessed by the
Orthodox and by the classical Lutheran confession, and is widely, though
not universally, held in the wider Reformation inheritance today; the
Immaculate Conception is rejected by the Ecumenical Patriarchate's 1895
encyclical as an innovation and questioned by Anglicans as to its
required assent; the Assumption's content is largely shared and its
definitional mode contested. The fault line, in short, is not Mariology
but the theology of development and the locus of teaching
authority---which is why LG 67 warns preachers to ``keep away from
whatever, either by word or deed, could lead separated brethren or any
other into error regarding the true doctrine of the Church,'' and why the
Church's current discipline on Marian titles (Chapter 8) is framed with
the same care.\endnote{LG 67, Vatican English text, same locus. The
characterization of present Reformation practice is a general
observation, not a documented survey, and is offered at that level.}

\subsection*{Currentness}

The one mutable element in this reference is the Holy See's present
judgment on Marian cooperation titles. As of the verification date, the
controlling document is the doctrinal note \work{Mater Populi Fidelis} of
the Dicastery for the Doctrine of the Faith, approved by Pope Leo XIV on
7 October 2025 and published 4 November 2025, treated in Chapter 8. Its
text, approval, publication date, and signatories were re-verified on the
Holy See's own page on the verification date, and a search for a later
Holy See act on Marian titles returned none; a reader consulting this
reference later should recheck, since silence is not evidence of
continuance.\endnote{\work{Mater Populi Fidelis}, re-read 2026-07-25 at
\url{https://www.vatican.va/roman_curia/congregations/cfaith/documents/rc_ddf_doc_20251104_mater-populi-fidelis_en.html}
(approval by Leo XIV, 7 October 2025; publication 4 November 2025;
signed by Cardinal V\'ictor Manuel Fern\'andez, Prefect, and Mons.
Armando Matteo, Secretary for the Doctrinal Section). The negative
currentness check was a web search for later Holy See acts on Marian
titles, which returned only the 2025 note and coverage of it; the Vatican
Press Office bulletin of 4 November 2025 announces the same document. The
Dicastery's general document index page returned an empty body on the
verification date---a checked negative that limits the strength of the
search, and is recorded rather than concealed. The four dogmas themselves
are stable doctrine and are not given artificial currentness.}
